\ncdlesson{15}{Nyheder i algoritmernes tidsalder: Ekkokamre og kildekritik}{算法时代的新闻：回音室与信源批判}
{DU3 Modul 4（B1+ $\rightarrow$ B2-）｜对标 NCE 3 \& Modultest 4 \textit{Læsning \& Skrivning Opgave 2（批判性思辨辩论）}}
{\textbf{批判性思维核心表达}（\textit{sætte spørgsmålstegn ved, forholde sig kritisk til, nuancere debatten}）· 媒体素养惯用语（\textit{tage for gode varer, skille fakta fra holdninger}）· Unit 3 收官综合思辨}
{欢迎来到 Unit 3“科技、媒体与现代伦理”的收官课！在算法主宰信息流的今天，人们极易被困在只放大同温层声音的“回音室（\textit{ekkokammer}）”中。在 DU3 Modul 4 考试中，如何用丹麦语冷静地提出质疑、区分事实与立场、并使讨论摆脱偏激的非黑即白？掌握本课的 \textbf{\textit{sætte spørgsmålstegn ved}}、\textbf{\textit{forholde sig kritisk til}} 与 \textbf{\textit{nuancere debatten}}，你就真正具备了丹麦语 B2- 公民思辨力！}

% ---------------------------------------------------------------------
\begin{foerduleaeser}
带着下面两个问题阅读 Tekst A 和 Tekst B：
\begin{enumerate}
  \item \textit{Hvordan opstår ekkokamre på sociale medier ifølge Tekst A, og hvorfor kan de skade den demokratiske samtale?}（根据 Tekst A，社交媒体上的“回音室”是如何产生的？为什么它们可能损害民主对话？）
  \item \textit{Hvilke tre konkrete vaner foreslår journalisten Sara i Tekst B, når man møder en opsigtsvækkende nyhed på nettet?}（在 Tekst B 中，记者萨拉建议大家在网上看到一条耸人听闻的新闻时，养成哪三个具体习惯？）
\end{enumerate}
\end{foerduleaeser}

% ---------------------------------------------------------------------
\ncdsubsec{二、Tekster（科技媒体专题与思辨评论双语篇）}

\subsection*{Tekst A：\textit{Nyheder i algoritmernes tidsalder: Ekkokamre og kildekritik}（媒体思辨专栏 · 约 255 词）}

\begin{tipbox}{语篇说明}
本篇聚焦算法推荐机制下的信息茧房与信源批判意识。阅读时请重点标出\textbf{\color{NCDNavy}深蓝色加粗的批判性思维核心动词短语}（留意介词 \textit{ved} 和反身代词 \textit{sig}），以及\textbf{\color{NCDRed}深红色加粗的习惯搭配与从句衔接}。
\end{tipbox}

For tyve år siden fik de fleste danskere deres daglige nyheder fra trykte aviser og faste tv-udsendelser. I dag orienterer især mange unge sig hurtigt gennem personlige feeds på sociale medier. Det giver nem og øjeblikkelig adgang til verdens begivenheder, men eksperter opfordrer borgerne til at \da{forholde sig kritisk til} den måde, nyhederne bliver udvalgt på.

Bag ethvert nyhedsfeed arbejder nemlig usynlige algoritmer. De registrerer, hvilke opslag vi klikker på, og viser os derefter endnu mere indhold, som ligner det samme. På den måde risikerer brugerne at havne i såkaldte ekkokamre, hvor deres egne fordomme hele tiden bliver bekræftet, uden at de møder modargumenter.

Når stærke følelser og hurtige delinger belønnes med flest klik, bliver det sværere at \dabold{skille fakta fra holdninger}. Derfor bør enhver læser \da{sætte spørgsmålstegn ved}, hvem der står bag en påstand, og hvad afsenderens formål egentlig er. Man må simpelthen aldrig \dabold{tage} en dramatisk overskrift \dabold{for gode varer}, før man har tjekket oplysningerne hos uafhængige medier.

Heldigvis kan man aktivt bryde ud af sin egen boble. Redaktionelle medier og bred public service-journalistik spiller fortsat en vigtig rolle, fordi uddannede journalister arbejder med at efterprøve kilder og \da{nuancere debatten}. Jo bedre vi træner vores kildekritik i skolen og i hverdagen, desto stærkere står det oplyste demokrati.

\begin{transbox}
\small
\textbf{\color{NCDTeal}[Tekst A 中文精译]}\\
二十年前，大多数丹麦人通过纸质报纸和固定时段的电视节目获取每日新闻。而如今，尤其是许多年轻人习惯通过社交媒体上的个性化信息流快速了解时事。这虽然提供了便捷即时的世界资讯入口，但专家呼吁公民对新闻被筛选的方式保持批判态度。

因为在每一条信息流的背后，都有看不见的算法在运转。算法会记录我们点击了哪些帖子，随后向我们推送更多相似的内容。如此一来，用户就可能落入所谓的“回音室”——在其中，他们自己的偏见不断得到印证，却很难接触到相反的论点。

当强烈情绪与快速转发能换来最多点击量时，将事实与观点区分开来就变得愈发困难。因此，每位读者都应当对某个论断背后是谁、发布者真正的意图是什么提出质疑。人们绝不应在尚未向独立媒体核查信息之前，就轻信一个耸动的标题。

幸好，人们可以主动打破自己的信息气泡。有专业采编团队的媒体以及覆盖面广的公共服务新闻依然扮演着重要角色，因为受过专业训练的记者致力于核实信源，并让公共讨论呈现出多维层次。我们在学校和日常生活中把信源批判力训练得越好，建立在充分知情基础上的民主制度就越稳固。
\end{transbox}

\subsection*{Tekst B：\textit{Sådan træner du din kildekritik}（媒体素养访谈节选 · 约 140 词）}

\begin{emailbox}{校园广播访谈｜Jonas（大学生）\& Sara（新闻记者）}
\textbf{Jonas:} Sara, mange unge scroller bare hurtigt forbi ti overskrifter i minuttet. Hvordan undgår man egentlig at blive snydt af falske nyheder?\\
\textbf{Sara:} Først og fremmest skal man lære at stoppe op og \da{sætte spørgsmålstegn ved}, om en historie lyder for utrolig til at være sand. Vær ekstra skeptisk, hvis en tekst gør dig vred med det samme.\\
\textbf{Jonas:} Hvad gør jeg så, hvis jeg gerne vil tjekke en påstand, før jeg deler den med vennerne?\\
\textbf{Sara:} Så bør du altid \da{forholde dig kritisk til} kilden og tjekke mindst to uafhængige nyhedsmedier. Derudover hjælper det at læse analyser, der forsøger at \da{nuancere debatten}, i stedet for kun at følge korte videoklip, der giver dig ret.
\end{emailbox}

\textbf{Tekst B 阅读核对（Rigtigt / Forkert / Kort svar，答案见第六部分）：}
\begin{enumerate}
  \item \textit{Sara anbefaler, at man straks deler nyheder, hvis de vækker stærke følelser eller vrede.} \underline{\hspace{2.2cm}}
  \item \textit{Hvor mange uafhængige nyhedsmedier bør man mindst tjekke ifølge Sara?} \underline{\hspace{2.2cm}}
  \item \textit{Hvorfor bør man vælge analyser frem for kun korte videoklip?} \underline{\hspace{2.2cm}}
\end{enumerate}

\begin{tipbox}{本课发音与语调要点（Udtale \& Reduktioner）}
\begin{itemize}
  \item \textbf{批判性思维核心词与软 \textit{d} / 不发音辅音}：
    \begin{itemize}
      \item \textit{et spørgsmålstegn} \ipa{ˈsbɶɐ̯smɔːlsˌtˢɑi̯ˀn}（\textit{spørgs-} 中的 \textit{g} 不发音；\textit{tegn} 读滑音 \ipa{tˢɑi̯ˀn}，\textit{g} 消失化作长音与喉塞音！）
      \item \textit{at nuancere} \ipa{nyɑŋˈseːˀɐ}（法语借词！\textit{nu-} 读轻柔的 \ipa{ny-}，\textit{-an-} 读鼻音 \ipa{-ɑŋ-}，重音在 \textit{-ce-}）
      \item \textit{en fordom} \ipa{ˈfɒːˌdʌmˀ}（重音在前面 \textit{for-}，不同于英文的后缀重音！）
      \item \textit{en påstand} \ipa{ˈpʰʌˌsdanˀ}（词尾 \textit{d} 不发音，带 \textit{stød}！）
    \end{itemize}
  \item \textbf{媒体与信息高频词}：\textit{en algoritme} \ipa{ɑlɡoˈʁedmə}；\textit{et ekkokammer} \ipa{ˈɛɡoˌkʰɑmˀɐ}；\textit{kildekritik} \ipa{ˈkʰeləˌkʰʁitˢiɡ}（第一个 \textit{d} 不发音！）；\textit{troværdig} \ipa{tˢʁoˈvɛɐ̯ˀdi}（词尾 \textit{g} 不发音）。
\end{itemize}
\end{tipbox}

\begin{transbox}
\small
\textbf{\color{NCDTeal}[Tekst B 中文对照]}\\
\textbf{乔纳斯}：萨拉，很多年轻人一分钟内就快速划过十来条新闻标题。到底怎样才能避免被假新闻欺骗呢？\\
\textbf{萨拉}：首先，人必须学会停下来并提出质疑：一个故事听起来是不是好得或离奇得不像真的。如果一段文字让你立刻感到愤怒，就要格外保持怀疑。\\
\textbf{乔纳斯}：那如果我想在转发给朋友之前核实某个说法，该怎么做呢？\\
\textbf{萨拉}：那你就应当始终对信息来源保持批判态度，并且至少核对两家独立的新闻媒体。此外，阅读那些试图让讨论更全面深入的分析文章会很有帮助，而不要只关注一味顺着你心意的短视频片段。
\end{transbox}

% ---------------------------------------------------------------------
\ncdsubsec{三、Gloser（核心生词表）}

\begin{tipbox}{词表重点}
注意词尾双写与变音规则：\textit{en fordom}（偏见）定指和复数要双写 \textit{m}（\textit{fordommen, fordomme, fordommene}）；\textit{at forholde sig til} 是反身动词短语，绝对不能漏掉反身代词 \textit{sig}；而 \textit{en kildekritik}（信源批判能力）是不可数共性名词。
\end{tipbox}

\begin{glosertable}
名词 & en algoritme & \textit{(algoritmen, algoritmer, algoritmerne)} & 算法 \\
名词 & et ekkokammer & \textit{(ekkokammeret, ekkokamre, -kamrene)} & 回音室，信息茧房 \\
名词 & en kildekritik & \textit{(kildekritikken, ingen flertal)} & 信源批判，信息甄别能力 \\
名词 & en kilde & \textit{(kilden, kilder, kilderne)} & 来源，信源，消息源 \\
名词 & et spørgsmålstegn & \textit{(spørgsmålstegnet, spørgsmålstegn, -ene)} & 问号，质疑（单复同形） \\
名词 & en fordom & \textit{(fordommen, fordomme, fordommene)} & 偏见，成见 \\
名词 & en påstand & \textit{(påstanden, påstande, påstandene)} & 论断，主张，说法 \\
名词 & en dækning & \textit{(dækningen, dækninger, dækningerne)} & 新闻报道，覆盖面 \\
名词 & en troværdighed & \textit{(troværdigheden, ingen flertal)} & 可信度，公信力 \\
名词 & en offentlighed & \textit{(offentligheden, ingen flertal)} & 公众舆论界，社会公众 \\
\midrule
动词 [har] & at nuancere & \textit{(nuancerer, nuancerede, har nuanceret)} & 使具层次感，多维度阐明 \\
动词 [har] & at forholde sig (til) & \textit{(forholder, forholdt, har forholdt sig til)} & 对……持……态度，应对 \\
动词 [har] & at bekræfte & \textit{(bekræfter, bekræftede, har bekræftet)} & 证实，印证，强化 \\
动词 [har] & at efterprøve & \textit{(efterprøver, efterprøvede, har efterprøvet)} & 核查，检验，复核 \\
动词 [har] & at skille & \textit{(skiller, skilte, har skilt)} & 区分，分离（\textit{skille A fra B}） \\
\midrule
形容词 & troværdig & \textit{(troværdigt, troværdige)} & 可信的，可靠的 \\
形容词 & ensidig & \textit{(ensidigt, ensidige)} & 片面的，单方面的 \\
\end{glosertable}

% ---------------------------------------------------------------------
\ncdsubsec{四、Grammatik \& Stilistik（核心语法、伦理表达与修辞精讲）}

\begin{grammatikbox}{语法点 1：批判性思维三大王牌句型与介词骨架}
进入 DU3 Modul 4 辩论与阅读深水区，你必须精准驾驭下面三组核心批判性思维动词短语。请把它们的\textbf{固定介词、反身成分与从句接口}背成肌肉记忆：

\vspace{0.3em}
\renewcommand{\arraystretch}{1.25}
{\small
\begin{tabularx}{\linewidth}{@{} >{\bfseries\raggedright\arraybackslash}p{3.3cm} >{\raggedright\arraybackslash}p{3.9cm} >{\raggedright\arraybackslash}X @{}}
\toprule
\textbf{核心表达} & \textbf{句法结构与接续规则} & \textbf{标准 B1+ / B2- 示范句} \\
\midrule
\da{at sætte spørgsmålstegn ved}\newline （对……提出质疑） &
后接名词，或直接接 \textbf{\textit{, om / hvorvidt}} 从句；介词固定为 \textbf{\textit{ved}} &
\textit{Man bør \da{sætte spørgsmålstegn ved} kildens formål.}\newline
\textit{Jeg \da{sætter spørgsmålstegn ved, om} tallene passer.} \\
\midrule
\da{at forholde sig kritisk til}\newline （对……保持批判态度） &
必须带反身代词（\textit{mig / dig / sig / os}）+ 介词 \textbf{\textit{til}} &
\textit{Unge skal lære at \da{forholde sig kritisk til} nyheder på nettet.}\newline
\textit{（可替换副词：\textit{forholde sig skeptisk / neutralt til}）} \\
\midrule
\da{at nuancere debatten}\newline （使讨论更具层次/避免片面） &
直接及物动词（不加介词！），常见宾语：\textit{debatten, billedet, diskussionen} &
\textit{Artiklen er med til at \da{nuancere debatten} om algoritmer.}\newline
（也可以用分词定语：\textit{en \da{nuanceret} dækning}） \\
\bottomrule
\end{tabularx}
}
\end{grammatikbox}

\begin{grammatikbox}{语法点 2：媒体批判常用习语与 Unit 3 思辨句法联动}
把本课批判思维表达与前面几课的武器结合起来，就能写出真正 B2- 水平的 \textit{Debatindlæg}：
\begin{enumerate}
  \item \textbf{\da{at tage noget for gode varer}（不加思考地轻信某事）}：极地道的丹麦语惯用语（多用于否定或警示句）：\\
    \textit{Vi må ikke bare \da{tage} opslag på sociale medier \da{for gode varer}.}
  \item \textbf{\da{at skille fakta fra holdninger}（把事实与观点剥离开来）}：\\
    \textit{En god kildekritik hjælper os med at \da{skille fakta fra holdninger}.}
  \item \textbf{联动 Unit 3 前几课句法升级论证}：\\
    \ding{172} 结合 L11 名词化：\textit{\dabold{Ved efterprøvelse af} kilderne undgår man at sprede usande påstande.}\\
    \ding{173} 结合 L12 前提条件：\textit{Sociale medier informerer hurtigt, \dabold{forudsat at} vi forholder os kritisk til indholdet.}
\end{enumerate}
\end{grammatikbox}

\begin{grammatikbox}{英丹/中丹对比避坑指南（Stilistik \& Pitfalls）}
\begin{enumerate}
  \item \textbf{\textit{sætte spørgsmålstegn} 后面只能接 \dabold{ved}}：英语说 \textit{put a question mark on / question something}，中文说“对……画问号”，丹麦语介词固定且唯一是 \dabold{ved}：\ding{55}\ \textit{*sætte spørgsmålstegn på / til påstanden} $\rightarrow$ \ding{51}\ \textit{sætte spørgsmålstegn \da{ved} påstanden}。后面接“是否……”时，必须用 \da{om} 或 \da{hvorvidt}，严禁写成 \textit{*ved hvis}！
  \item \textbf{千万别漏掉反身代词 \dabold{sig}（或 \textit{mig / os}）}：丹麦语中“持某种态度面对某事”是反身结构 \textit{at forholde \dabold{sig} [副词] til}：\ding{55}\ \textit{*Jeg forholder kritisk til medierne} $\rightarrow$ \ding{51}\ \textit{Jeg forholder \dabold{mig} kritisk til medierne}。注意此处的 \textit{kritisk} 作副词修饰动作，因此遇到复数主语时（如 \textit{Vi forholder os \dabold{kritisk} til...}）通常保持基本形 \textit{kritisk}（不加 \textit{-e}）！
  \item \textbf{\textit{nuancere} 不是“中立/和稀泥”}：它表示“展现事物的复杂性与多面性，避免单向偏激”，相当于英文的 \textit{add nuance to}，而非“削弱立场”。
\end{enumerate}
\end{grammatikbox}

% ---------------------------------------------------------------------
\ncdsubsec{五、Mønsterøvelser（句型与修辞操练）}

\textbf{练习 A：介词与批判性思维核心词选填（从 \textit{ved, til, sig, fra, for, nuancere, fordomme, efterprøve} 中选出最恰当的词）}
\begin{enumerate}
  \item Læserne bør altid sætte spørgsmålstegn \underline{\hspace{2.2cm}} anonyme påstande på sociale medier.
  \item Som borger i et demokrati skal man lære at forholde \underline{\hspace{2.2cm}} kritisk \underline{\hspace{2.2cm}} nyhedsdækningen.
  \item En grundig baggrundsartikel kan være med til at \underline{\hspace{2.2cm}} debatten om kunstig intelligens.
  \item Man må aldrig tage dramatiske overskrifter \underline{\hspace{2.2cm}} gode varer uden at tjekke kilden.
  \item Algoritmer viser os ofte historier, der blot bekræfter vores egne \underline{\hspace{2.2cm}}.
  \item Kildekritik handler blandt andet om at skille dokumenterede fakta \underline{\hspace{2.2cm}} personlige holdninger.
  \item Dygtige journalister bruger tid på at \underline{\hspace{2.2cm}} oplysninger hos flere uafhængige kilder.
  \item Jeg sætter spørgsmålstegn \underline{\hspace{2.2cm}}, hvorvidt korte videoklip kan erstatte rigtige aviser.
\end{enumerate}

\textbf{练习 B：思辨表达升级手术（用本课批判性思维表达改写下列口语化简单句，或修正语法错误）}
\begin{enumerate}
  \item 用 \textit{sætte spørgsmålstegn ved, om...} 升级：\textit{Jeg tror ikke på, at denne video viser hele sandheden.} $\rightarrow$
  \item 用 \textit{forholde sig kritisk til} 升级：\textit{Unge bør ikke bare tro på alt, hvad de ser på deres feed.} $\rightarrow$
  \item 用 \textit{nuancere debatten} 升级：\textit{Vi har brug for artikler, som viser både fordele og ulemper i diskussionen.} $\rightarrow$
  \item 用 \textit{tage noget for gode varer} 改写：\textit{Du må ikke acceptere en overskrift uden at tjekke den først.} $\rightarrow$
  \item 纠错改写（漏了反身代词与介词错误）：\textit{Mange borgere forholder kritisk om de hurtige nyheder i ekkokamre.} $\rightarrow$
  \item 纠错改写（介词错误与误用 \textit{hvis}）：\textit{Eksperten sætter spørgsmålstegn på, hvis algoritmer kan være helt neutrale.} $\rightarrow$
\end{enumerate}

\textbf{练习 C：中译丹（综合调用批判性思维表达与新闻词汇）}
\begin{enumerate}
  \item 我们应当始终对网上来历不明的说法提出质疑。（用 \textit{sætte spørgsmålstegn ved}）
  \item 年轻人需要学会对社交媒体上的新闻保持批判态度。（用 \textit{forholde sig kritisk til}）
  \item 独立的报道有助于让公共讨论变得更有层次。（用 \textit{nuancere debatten}）
  \item 在核查信源之前，绝不能轻信耸人听闻的标题。（用 \textit{tage ... for gode varer}）
  \item 算法往往会强化我们自己的偏见，并制造信息茧房。
  \item 学会区分事实与个人观点是现代公民的重要能力。
\end{enumerate}

% ---------------------------------------------------------------------
\ncdsubsec{六、DU3 Modul 4 实战微任务与参考答案}

\begin{emailbox}{DU3 Modul 4 微型写作任务（对标 Skrivning Opgave 2 · Debatindlæg）}
\textbf{Situation}: På et ungdomsmagasins hjemmeside skriver bloggeren Freja:\\
\textit{”Jeg får alle mine nyheder på fem minutter gennem korte videoer på sociale medier. Traditionelle aviser og tv-nyheder er langsomme, gammeldags og helt unødvendige i dag.”}\\[4pt]
\textbf{Opgave}: Skriv et debatindlæg på \textbf{mindst 100 ord} (anbefalet 130--160 ord), hvor du tager stilling til Frejas påstand. Forklar både muligheder og farer ved nyheder på sociale medier, og brug mindst \textbf{tre} udtryk for kritisk tænkning fra denne lektion (\textit{sætte spørgsmålstegn ved, forholde sig kritisk til, nuancere debatten, skille fakta fra holdninger, tage noget for gode varer}).
\end{emailbox}

\begin{answerbox}
\textbf{Før du læser 参考回答}
\begin{enumerate}
  \item \textit{Ekkokamre opstår, fordi algoritmer registrerer vores klik og derefter viser mere indhold, som bekræfter vores egne fordomme. Det skader demokratiet, fordi borgerne ikke længere møder modargumenter og har sværere ved at skille fakta fra holdninger.}
  \item \textit{Sara foreslår, at man stopper op og sætter spørgsmålstegn ved utrolige historier, tjekker mindst to uafhængige nyhedsmedier samt læser analyser, der nuancerer debatten.}
\end{enumerate}

\textbf{Tekst B 阅读核对答案}
\begin{enumerate}
  \item \da{Forkert}（相反，当一条文字让人立刻感到愤怒时，她建议人要格外保持怀疑）
  \item \da{Mindst to uafhængige nyhedsmedier.}（至少核查两家独立的新闻媒体）
  \item \da{Fordi analyser forsøger at nuancere debatten.}（因为分析文章试图让讨论更具层次感，而不是只一味附和你的观点）
\end{enumerate}

\textbf{练习 A 答案}
\begin{enumerate}
  \item \da{ved}（固定介词：\textit{sætte spørgsmålstegn ved}）
  \item \da{sig}；\da{til}（反身代词 \textit{sig} 呼应泛指主语 \textit{man}，介词固定为 \textit{til}）
  \item \da{nuancere}（\textit{at nuancere debatten} 使讨论更全面深入）
  \item \da{for}（习语：\textit{tage noget for gode varer} 轻信）
  \item \da{fordomme}（复数形式双写 \textit{m}：\textit{egne fordomme}）
  \item \da{fra}（固定搭配：\textit{skille A fra B} 把 A 从 B 中区分开）
  \item \da{efterprøve}（\textit{at efterprøve oplysninger} 核查核实信息）
  \item \da{ved}（接 \textit{hvorvidt} 从句前仍必须保留介词 \textit{ved}！）
\end{enumerate}

\textbf{练习 B 答案（句型升级参考答案）}
\begin{enumerate}
  \item \textit{Jeg \da{sætter spørgsmålstegn ved, om} denne video viser hele sandheden.}（用 \textit{ved, om...} 升级客观质疑）
  \item \textit{Unge bør \da{forholde sig kritisk til} alt, hvad de ser på deres feed.}（准确使用反身代词 \textit{sig} 与介词 \textit{til}）
  \item \textit{Vi har brug for artikler, som \da{nuancerer debatten}.}（用凝练动词 \textit{nuancere} 替代松散描述）
  \item \textit{Du må ikke \da{tage} en overskrift \da{for gode varer} uden at tjekke den først.}（地道习语替换）
  \item \textit{Mange borgere \dabold{forholder sig} kritisk \dabold{til} de hurtige nyheder i ekkokamre.}（补齐反身代词 \textit{sig}，并将错介词 \textit{om} 改为 \textit{til}）
  \item \textit{Eksperten sætter spørgsmålstegn \dabold{ved, hvorvidt} algoritmer kan være helt neutrale.}（介词 \textit{ved} + 表达“是否”必须用 \textit{hvorvidt} 或 \textit{om}，不可用 \textit{hvis}）
\end{enumerate}

\textbf{练习 C 答案}
\begin{enumerate}
  \item \textit{Vi bør altid sætte spørgsmålstegn ved anonyme påstande på nettet.}
  \item \textit{Unge skal lære at forholde sig kritisk til nyheder på sociale medier.}
  \item \textit{Uafhængig journalistik hjælper med at nuancere den offentlige debat.}
  \item \textit{Man må aldrig tage dramatiske overskrifter for gode varer før kildekritik.}
  \item \textit{Algoritmer bekræfter ofte vores egne fordomme og skaber lukkede ekkokamre.}
  \item \textit{At skille fakta fra personlige holdninger er en vigtig evne for borgere.}
\end{enumerate}

\textbf{DU3 Modul 4 高分范文（Debatindlæg · 150 词）}\\
\textit{\textbf{Hurtige nyheder gør os ikke klogere}}

\textit{Freja har jo ret i, at sociale medier giver et hurtigt overblik over dagens begivenheder, og mange unge foretrækker korte videoer frem for lange artikler i en travl hverdag.}

\textit{Alligevel mener jeg, at traditionelle nyhedsmedier stadig er helt uundværlige. På sociale medier styrer algoritmer nemlig, hvad vi ser, og vi risikerer let at havne i ekkokamre, der kun bekræfter vores egne fordomme. Derfor bør man altid \da{forholde sig kritisk til} korte klip og \da{sætte spørgsmålstegn ved}, om de overhovedet fortæller hele sandheden.}

\textit{Uddannede journalister arbejder med at efterprøve kilder og \da{skille fakta fra holdninger}. Hvis vi stopper med at læse dybdegående journalistik, mister vi de analyser, der kan \da{nuancere debatten}. Sociale medier kan altså supplere, men aldrig erstatte troværdige aviser.}

\small\textbf{考官点评}：全篇 150 词（不含标题），完美展示 Unit 3 收官思辨水准！开篇先用情态小品词 \textit{jo} 认可社交媒体“快速概览”的便利，继而指出算法与回音室隐患，连续精准调度 \textit{forholde sig kritisk til}、\textit{sætte spørgsmålstegn ved, om...}、\textit{skille fakta fra holdninger} 与 \textit{nuancere debatten} 四大武器，论证逻辑层层递进、收束利落有力。
\end{answerbox}
